\section{9.1 Thesis Summary}
This thesis built and validated a full four-camera perception pipeline for body-part-targeted football training research. The work integrated the full \texttt{Project\_Cam} history into one coherent framework and delivered measurable performance under real garage constraints.

Main outputs include:
\begin{itemize}
  \item robust calibration workflows,
  \item synchronized 4-camera data acquisition,
  \item 3D ball and skeleton triangulation,
  \item arena-aware rendering and diagnostics,
  \item GT protocols and quantitative reports.
\end{itemize}

\section{9.2 Scientific Contributions}
\begin{enumerate}
  \item A repository-wide integration narrative from prototype to validated pipeline.
  \item A practical calibration methodology with robust extrinsics and partial camera recalibration support.
  \item A protocol-driven GT framework for both rigid-like ball tests and joint-touch tests.
  \item Quantitative evidence distinguishing raw performance, bias-corrected performance, and dynamic stress behavior.
  \item A concrete, safety-aware roadmap to launcher integration.
\end{enumerate}

\section{9.3 Limitations}
\begin{itemize}
  \item One camera remained more sensitive to drift and limited tag overlap.
  \item Dynamic fast ball throws still produce outliers.
  \item Joint GT includes unavoidable human placement uncertainty.
  \item Closed-loop launcher control was not fully implemented in this thesis window.
\end{itemize}

These limitations are explicitly documented to avoid over-claiming.

\section{9.4 Future Work}
\subsection{Immediate}
\begin{itemize}
  \item complete rigid-point campaign after final mount lock,
  \item improve overlap for >=3-camera support across operational volume,
  \item finalize south-camera stabilized recalibration procedure,
  \item refine dynamic outlier handling with physically informed filtering.
\end{itemize}

\subsection{Mid-Term}
\begin{itemize}
  \item implement command parser and target semantic mapper,
  \item calibrate launcher frame and release model,
  \item run HIL supervised firing tests with visual verification.
\end{itemize}

\subsection{Long-Term}
\begin{itemize}
  \item adaptive user-specific targeting profiles,
  \item predictive motion compensation,
  \item curriculum-based autonomous training sessions.
\end{itemize}

\section{9.5 Final Statement}
The project has reached a meaningful milestone: perception and evaluation are now organized, quantified, and reproducible. This is the necessary scientific and engineering foundation for safe intelligent actuation. The remaining work is substantial but clearly defined, making the transition from research prototype to functional smart training system realistic.


\section{9.6 Final Research Positioning}
The thesis should be interpreted as a completed perception-and-validation milestone in a larger intelligent training system program. The foundational question "Can we build and measure a deployable 3D perception stack in this environment?" is answered positively with quantified limitations.

\section{9.7 What This Enables Immediately}
Immediate enabled capabilities include:
\begin{itemize}
  \item software-defined target-zone analytics,
  \item objective player-session review in 3D,
  \item controlled experiments for command-driven aiming logic,
  \item robust baseline for publishing system-integration findings.
\end{itemize}

\section{9.8 Closing Remark}
A strong research thesis in engineering is not the one that claims everything is solved; it is the one that clearly defines what is solved, what is measured, what remains uncertain, and what the next technically sound step is. This manuscript follows that principle.



\section{9.9 Final Outcome Statement for Stakeholder and Academic Audiences}
For stakeholders, the project now provides a practical and measurable 3D perception platform that can support smarter training workflows. For academic evaluation, the thesis provides a complete research cycle: hypothesis, implementation, protocolized experimentation, quantitative analysis, and a technically grounded roadmap.

The final claim is therefore precise: the perception backbone is validated and integration-ready, while full autonomous launcher control is the next phase, not a completed component of this manuscript.
